% =====================================================================
%  Results -- Characterisation of a large-scale RT/retron mining corpus
%
%  Drop-in for a thesis Results chapter. Written at \section level:
%  if this is its own chapter, change \section->\chapter and demote the rest.
%
%  Requires:
%    \usepackage{booktabs}
%    \usepackage{graphicx}
%    \usepackage{siunitx}          % for \num{} and \SI{}{\percent}
%    \sisetup{group-separator={,}} % thousands separator
%  siunitx v3 users: add [=v2] i.e. \usepackage[=v2]{siunitx}, or replace
%  \SI{x}{\percent} with \qty{x}{\percent} throughout.
%  If you prefer no siunitx: \newcommand{\num}[1]{#1} and
%  \newcommand{\SI}[2]{#1\%} reproduce the text with manual separators.
%  Figure files are referenced by name only; set \graphicspath to the
%  directory holding the exported figures, e.g.
%    \graphicspath{{figures/}}
%  Filenames contain underscores. This is fine in modern LaTeX; if your
%  setup objects, add \usepackage{grffile} or rename the files.
%
%  Every number here is reproduced by the analysis notebook; the
%  figure/table manifest lists the source table for each.
% =====================================================================

\section{Characterisation of the reverse-transcriptase mining corpus}
\label{sec:dbchar}

\subsection{Corpus definition and analytical populations}
\label{sec:dbchar:corpus}

The corpus comprises 43 newline-delimited JSON files totalling \num{81007695609} bytes and
\num{3358182} records, obtained by mining genomes and metagenome-assembled genomes from eight
source collections: \texttt{ncbi\_bacteria}, \texttt{gtdb\_bacteria}, \texttt{ncbi\_archaea},
\texttt{gtdb\_archaea}, \texttt{mgnify\_human\_gut}, \texttt{mgnify\_soil},
\texttt{mgnify\_marine} and \texttt{gem}. Each file was pinned by its SHA-256 digest and the
record set by a single digest computed over the file, line number, byte offset and content hash of
every record, so that all downstream quantities are attributable to a fixed input state. No record
failed to parse.

The corpus separates into three file-pure anchor populations (Table~\ref{tab:populations}):
\num{3050688} RT-anchored records distributed across 41 single-family files, \num{9012} RT-anchored
records carrying multiple family labels, and \num{298482} records anchored on a non-coding RNA
rather than on a reverse transcriptase.

The ncRNA-anchored records were excluded from every RT denominator reported in this chapter. They
are retained with an explicit population flag and may be analysed under their own denominator, but
they do not contribute to any RT rate, family statistic or RT--ncRNA geometry. Multi-labelled
records were likewise held out as a separate stratum (Section~\ref{sec:dbchar:multi}) and were never
merged into a single family.

Analyses requiring genomic context were computed on a working population defined by three
conditions: retention of one copy of each byte-identical duplicated record, exclusion of
multi-labelled records, and verification of the RT coordinates against the extracted window. This
population retains \num{3026000} records, corresponding to \num{2822547} loci, \num{2451078}
physical loci, \num{477956} distinct RT proteins and \num{1535827} genomes, and costs \num{33700}
records, or \SI{1.10}{\percent} of the corpus. Analyses requiring only the protein sequence were
computed on the wider population that omits the coordinate-verification condition, retaining
\num{493964} distinct RT proteins; the difference of \num{16008} proteins is discussed in
Section~\ref{sec:dbchar:integrity}.

\subsection{Analytical units do not convert by a constant factor}
\label{sec:dbchar:units}

Because a single genomic position may be deposited under several database accessions, and a single
RT protein may occur at many positions, the corpus supports several counting units that are not
interchangeable (Table~\ref{tab:units}). The \num{3059700} raw records reduce to \num{3051238}
distinct records, \num{2847312} accession-defined loci, \num{2475684} physical loci and
\num{501561} distinct RT proteins, across \num{1542433} genomes and \num{2737189} RT taxonomic
occurrences.

The corpus therefore contracts approximately sixfold between records and distinct proteins, but this
factor is an artefact of aggregation rather than a conversion constant. Computed separately for each
source collection, the number of accession-defined loci per distinct RT protein ranges from 1.11 in
\texttt{mgnify\_soil} to 5.49 in \texttt{ncbi\_bacteria} (Figure~\ref{fig:units}), a fivefold spread
that tracks sequencing depth rather than any property of the enzymes. A rate reported on records is
consequently weighted by isolate-genome sequencing effort, whereas the same rate reported on
distinct proteins is not.

Genome and contig identifiers are not unique across source collections. GTDB re-publishes assemblies
obtained from NCBI, and RefSeq contig accessions carry an \texttt{NZ\_} prefix absent from the
corresponding GenBank record, so a locus defined on the contig accession counts one physical
position twice when it has been deposited under both schemes. All \num{371628} RefSeq/GenBank twin
pairs were confirmed by identical window DNA, identical RT protein sequence, identical window
coordinates and identical strand, with no disagreements, so the collapse to physical loci is
supported by evidence rather than assumed. No locus carries conflicting RT sequences, while
\num{203921} loci appear under more than one source collection.

That overlap is structurally simple: essentially all of it consists of the single pair
\texttt{ncbi\_bacteria} and \texttt{gtdb\_bacteria} (\num{200914} loci in the working population),
with \num{111394} genomes appearing under two collections. Per-collection counts are therefore
reported as coverage rather than as shares. Where a partition is required it is taken at the locus
or genome level under a stated primary-source rule; it is not taken at the protein level at all,
because \SI{45.23}{\percent} of distinct RT proteins occur in two or more collections
(Figure~\ref{fig:dboverlap}).

\subsection{Sequence integrity and the classification of unresolved records}
\label{sec:dbchar:integrity}

Every RT interval was independently re-derived by translating the stored window DNA and comparing
the product with the stored protein. Of the single-family RT-anchored records, \num{2532006}
back-translate to an exact match and only 618 mismatch; \num{3034472} of \num{3050688}
(\SI{99.47}{\percent}) have coordinates verified against a self-consistent window.

Records that cannot support a given measurement were classified rather than discarded. A total of
\num{31504} RT-anchored records carry no annotated RT coding sequence. Of these, \num{14188} were
recovered by back-translation and retain full genomic context, \num{16688} retain a usable protein
sequence but no defensible genomic context, and 628 remain ill-posed
(Table~\ref{tab:recovery}). Recovery was validated against an independent translation
implementation, which agreed on \num{1999} of \num{2000} sampled records, with no case resolved only
by a frame shift.

The sequence-only category is entirely attributable to metagenomic assembly fragmentation. All
\num{16688} such records originate from metagenome-derived collections
(\texttt{mgnify\_human\_gut}, $n=\num{15364}$; \texttt{mgnify\_marine}, $n=766$;
\texttt{mgnify\_soil}, $n=533$; \texttt{gem}, $n=25$) and none from isolate-derived collections,
where the corresponding rate is zero. Within the human-gut MAG collection they represent
\SI{10.81}{\percent} of records. The three underlying representation classes---an RT lying beyond
the retrieved contig ($n=\num{9128}$), an RT crossing a contig boundary ($n=\num{3636}$), and an
inverted window from which no sequence could be recovered ($n=\num{3924}$)---are all expected
consequences of short, fragmented assemblies on which a gene may extend past the available sequence.
These records are excluded only from analyses requiring genomic context and are retained wherever
the protein sequence alone is the unit of analysis.

It should be noted that verification here establishes agreement between two fields of a single
record. Where the upstream pipeline derived both from the same erroneous source they would agree and
remain incorrect; the procedure bounds internal inconsistency, not upstream correctness.

\subsection{Composition of the RT family landscape}
\label{sec:dbchar:families}

Family composition was computed on distinct RT proteins, since the record-level distribution
principally reports which organisms have been sequenced most intensively. The \num{501561} distinct
RT proteins comprise \num{493956} single-family proteins distributed across 41 family labels,
\num{7593} multi-labelled proteins, and 12 proteins observed under more than one family label in
single-family files.

Three families dominate protein-level diversity: RVT-GII (\SI{51.95}{\percent}), Retron
(\SI{15.85}{\percent}) and RVT-DGRs (\SI{15.41}{\percent}), together accounting for approximately
\SI{83}{\percent} of single-family proteins (Figure~\ref{fig:families}). The record-level and
protein-level distributions differ for every family, and the direction of that difference identifies
those families concentrated in repeatedly deposited isolate genomes.

RT protein length differs between families but does not separate them. The overall median is
385~amino acids (range 20--\num{10449}), with Retron at 342 and RVT-GII at 407, while several
ungrouped families occupy markedly longer regimes (RVT-UG5, 892; RVT-UG8, 606) consistent with
domain fusion (Figure~\ref{fig:lengths}). The interquartile ranges of the three largest families
overlap substantially, so length alone does not constitute a classifier. Extreme values were
retained rather than trimmed.

Length is confounded with call completeness. Completeness is derived from the Prodigal
\texttt{partial} attribute, a two-digit code in which the first digit describes the 5$'$ end of the
called gene and the second the 3$'$ end; \texttt{0} indicates that the boundary was identified and
\texttt{1} that the gene runs off the edge of the contig. Thus \texttt{00} denotes a complete open
reading frame, \texttt{10} and \texttt{01} genes truncated at the 5$'$ and 3$'$ end respectively,
and \texttt{11} a gene whose contig is shorter than the gene itself. Across distinct records,
\SI{88.19}{\percent} are \texttt{00}, \SI{4.79}{\percent} \texttt{10}, \SI{4.31}{\percent}
\texttt{01} and \SI{1.68}{\percent} \texttt{11}. Over single-family proteins, \SI{25.2}{\percent}
are partial at every occurrence, and this fraction varies strongly between families, tracking those
whose length medians are lowest (Figure~\ref{fig:completeness}). Partiality is a property of the
assembly rather than of the protein, and a record lacking the attribute is classified as having no
completeness evidence rather than as complete. Family-level length comparisons intended as
statements about protein architecture should therefore be repeated on complete calls alone.

Thirty-six of 44 family length baselines established by earlier work were reproduced exactly,
providing an external consistency check on family assignment.

The twelve proteins observed under more than one family label all carry the same pair of labels,
\texttt{RVT-CRISPR} and \texttt{RVT-CRISPR-like}, indicating a single region of overlap between two
closely related profiles rather than dispersed labelling inconsistency. They range from 202 to 983
amino acids and occur in 2--50 genomes each.

\subsection{Multi-labelled proteins represent profile ambiguity}
\label{sec:dbchar:multi}

The \num{7593} multi-labelled distinct RT proteins were re-scored against the family profile
library. For \SI{99.91}{\percent} of them the highest-scoring family was among the labels already
assigned to the record, indicating that multiple labels reflect genuine profile ambiguity rather
than erroneous assignment.

The discriminating evidence is nonetheless weak. The median margin between the best- and
second-best-scoring family is 5~bits, against 81.7~bits in a seeded single-family control in which
the best-scoring profile matched the known label for \SI{98.10}{\percent} of proteins
(Figure~\ref{fig:multi}). Assigning each multi-labelled protein to its highest-scoring family would
therefore amount to an arbitrary choice presented as a classification.

Multi-labelled proteins are also shorter than single-family proteins (median 283 versus 385 amino
acids), and profile bit scores scale with alignment length, so the reduced margin is not fully
separable from a length effect; a length-matched control has not been performed. These proteins are
consequently excluded from family-level statistics, in which they would contribute to two
denominators simultaneously, and retained in sequence-level analyses where family membership is not
the unit of analysis.

\subsection{The non-coding RNA detection landscape}
\label{sec:dbchar:ncrna}

Retron non-coding RNAs were detected using a library of 21 covariance models, yielding \num{346722}
calls that reduce to \num{16458} distinct RNA sequences in the paired view. Every distinct RNA
sequence is matched by exactly one model, so per-model counts partition the sequence set. Called
RNAs occupy a narrow length regime (median 151~nt, interquartile range 131--193, range 34--395), and
most individual models produce an interquartile range narrower than approximately 20~nt
(Figure~\ref{fig:ncrnalen}). This is substantially a property of what a covariance model can match:
these are model-defined hit boundaries, not independently established transcript boundaries.
Structural annotation is present for only \SI{1.37}{\percent} of calls.

Call-level and sequence-level model composition diverge for the largest models, which produce many
more calls per distinct sequence than others; their matches are concentrated in repeatedly deposited
organisms and accumulate calls without accumulating diversity (Figure~\ref{fig:cmcomp}).

Detection is strongly family-restricted. Of \num{632688} Retron loci, \num{333838}
(\SI{52.77}{\percent}) carry at least one RNA call. Every other family lies at or above
\SI{99.9}{\percent} with no call (RVT-GII, \SI{99.99}{\percent}; RVT-DGRs, \SI{99.98}{\percent};
RVT-UG11, \SI{100.00}{\percent}). Among loci carrying a call, the large majority carry exactly one;
the zero and non-zero categories were verified to be mutually exclusive after aggregating all
records of each locus.

This contrast must not be read as biology. The covariance models are retron models, and outside the
Retron family the absence of a call measures the scope of the detector rather than the absence of an
associated RNA. Claims of biological absence in this corpus require a positive control demonstrating
that the same instrument recovers a known-present case on an appropriate substrate.

\subsection{RT--ncRNA genomic architecture}
\label{sec:dbchar:geometry}

All geometry was computed from coordinates and strand. The positional field shipped with the corpus
was not used, for reasons given in Section~\ref{sec:dbchar:technical}.

The \num{346722} placements reduce to \num{344154} canonical placements after geometry eligibility
and de-duplication. The dominant arrangement is close, same-strand and 5$'$ (Table~\ref{tab:geometry}):
the RNA lies upstream of the RT in \num{325269} placements (\SI{94.51}{\percent}) against \num{6761}
downstream and \num{12124} overlapping, at a median upstream separation of $-55$~bp, on the same
strand in \SI{99.12}{\percent} of cases, and with no annotated coding sequence between the two
features in \SI{94.41}{\percent} (Figure~\ref{fig:geometry}).

These four properties are strongly coupled rather than independently distributed. A single joint
combination---upstream, same strand, no intervening coding sequence, within 200~bp---accounts for
\SI{55.22}{\percent} of canonical placements, and the remaining combinations fall away by orders of
magnitude (Figure~\ref{fig:joint}).

The arrangement is not an artefact of repeated deposition. Re-weighting from placements to distinct
physical loci changes the marginals only marginally (upstream, \SI{94.51}{\percent} to
\SI{94.91}{\percent}; same strand, \SI{99.12}{\percent} to \SI{99.45}{\percent}). Re-weighting to
distinct sequence pairs is more discriminating and informative: the band within 200~bp strengthens
from \SI{55.22}{\percent} of placements to \SI{73.86}{\percent} of distinct pairs, whereas the
1--5~kb upstream band collapses from \SI{32.83}{\percent} to \SI{2.95}{\percent}, revealing that the
latter is carried by a small number of sequence pairs observed very many times. Close 5$'$
same-strand adjacency with no intervening coding sequence is thus the dominant arrangement under
every weighting, and the only band that strengthens when redundancy is removed.

These measurements describe genomic arrangement. They are compatible with, but do not demonstrate,
co-transcription, and they carry no information about cognate recognition between a given RT and a
given RNA.

Atypical arrangements were retained and quantified rather than removed. Within the eligible
placement population ($n=\num{345313}$), \num{12259} placements overlap the RT, \num{5903} contain
more than two intervening coding sequences, \num{3095} lie on the opposite strand and 557 are
separated by more than 5~kb (Table~\ref{tab:atypical}). These are heterogeneous in origin, including
window truncation, absent annotation in the intervening region, and possible biology; the
characterisation flags them rather than adjudicating among these causes.

\subsection{Two positional signals of technical origin}
\label{sec:dbchar:technical}

Two features of the data would yield confident but meaningless results if taken at face value.

First, the downstream separation distribution contains a narrow mode that is an artefact of window
extraction. All canonical downstream placements number \num{6761} and involve 230 distinct RNA
sequences. The mode itself, defined as the placements lying within 150~bp of the downstream median
of \num{2682}~bp, comprises \num{5234} placements (\SI{77.4}{\percent} of downstream placements)
involving only 18 distinct RNA sequences, 126 distinct RT proteins and 15 species, and is dominated
by a single stratum that is \SI{99.96}{\percent} truncated at the contig start. Where the extraction
window begins at the contig start, any call it contains is necessarily reported downstream of the
RT. The two quantities must be distinguished: the figure of 18 sequences characterises the mode, not
downstream placements in general.

Stratifying by the truncation flag also reveals a population that the artefact had obscured: the
\num{1041} downstream placements that are not truncated at the contig start lie at a median of only
64~bp and involve 149 distinct RNA sequences, a near-range downstream group that the technical mode
does not explain (Figure~\ref{fig:geometry}).

Second, the positional field shipped with the corpus is null for \num{331897} of \num{346722}
placements and corresponds to no coordinate frame; of twelve candidate frames tested, the closest
reproduces 146 values. Its variation decomposes as the negative offset of the RT within the window
plus the index of the intergenic region containing the RNA, that is, an index added to a coordinate.
It is retained as provenance only.

\subsection{Structure of the opposite-strand placements}
\label{sec:dbchar:opposite}

The \num{3041} canonical placements on the strand opposite to the RT are not randomly distributed.
\SI{84.5}{\percent} derive from a single covariance model, \texttt{TypeXIIIA\_firmi}, and
\SI{96}{\percent} occur at Retron loci. Their median separation from the RT is $-4750$~bp against
$-44$~bp for same-strand placements, and only \SI{12.5}{\percent} lack an intervening coding
sequence against \SI{95.1}{\percent} of same-strand placements. They involve 288 distinct RNA
sequences and 439 distinct RT proteins.

They are not a consequence of repeated deposition. The \num{2571} placements attributed to
\texttt{TypeXIIIA\_firmi} span \num{1324} distinct physical loci, \num{1322} genomes and 67 species
across two source collections, so the arrangement recurs across independent assemblies. More
strikingly, the same model additionally produces \num{1576} same-strand placements at a median
separation of 64~bp across 292 species, disjoint from the opposite-strand set in both separation and
taxonomic range (Figure~\ref{fig:strand}). A single covariance model therefore recovers two distinct
populations.

Whether the opposite-strand population reflects a genuine distant antisense arrangement restricted
to a narrow clade, or the model matching a second and unrelated element, cannot be resolved from
placement geometry alone and requires characterisation of the RNA sequences involved.

\subsection{Retron covariance-model hits at non-Retron RT loci}
\label{sec:dbchar:nonretron}

Non-Retron RT loci carrying a retron covariance-model hit show a weaker positional association:
\SI{57.69}{\percent} of 260 eligible single-family non-Retron placements lie upstream of the RT,
compared with \SI{94.51}{\percent} among Retron placements. The separation is also an order of
magnitude larger. Upstream non-Retron placements lie at a median of \num{1802}~bp from the RT and
downstream placements at \num{4281}~bp, against a median of 55~bp for Retron placements.

These hits are therefore not architecturally consistent with retron loci, and proximity alone does
not support their interpretation as divergent retrons. They are retained as candidate associations
for annotation follow-up. Establishing anything stronger would require covariance-model hit quality,
phylogenetic placement of the RT, and inspection of the intervening region.

\subsection{Pairing topology and the independence of recurrence}
\label{sec:dbchar:topology}

Reducing placements to distinct sequence pairs yields \num{30924} exact RT--RNA pairs involving
\num{29192} distinct RT proteins and \num{16458} distinct RNA sequences. This view is derived from
the eligible placement population; \num{30427} pairs are derivable from canonical placements alone,
the remaining 497 being represented only by placements removed during de-duplication.

The association is predominantly one-to-one as observed. Of the RT proteins in this view,
\num{28271} (\SI{96.85}{\percent}) are associated with exactly one RNA sequence and 921 with more
than one, with a maximum of 176. Conversely \num{13542} RNA sequences (\SI{82.28}{\percent}) are
associated with exactly one RT protein and \num{2916} with more than one, with a maximum of 705.
Decomposing the bipartite association graph into \num{14918} connected components yields \num{12079}
strictly one-to-one components, \num{2293} of many-to-one shape, 179 of one-to-many shape and 367 of
many-to-many shape. Ambiguity is thus concentrated rather than dispersed: the one-to-one components
constitute \SI{80.97}{\percent} of components but contain only \SI{41.38}{\percent} of the RT
proteins in the view (Figure~\ref{fig:topology}).

Both forms of multiplicity have identifiable and largely technical origins. For RT proteins with
more than one RNA partner, essentially every placement occurs at a locus carrying exactly one
distinct RNA sequence (\num{237965} placements, against 376 at loci carrying two or three), so the
multiplicity arises from the same protein occurring at many loci whose RNA calls differ slightly,
not from competing matches recorded at one locus. For 659 of the 921 such proteins
(\SI{71.6}{\percent}) the partner sequences span 5~nt or less in length, that is, they are boundary
variants of one RNA; only 110 have partners differing by more than 50~nt. The extreme case of 176
partners is a near-continuous length series from 68 to 223~nt with 57 sequences at exactly 158~nt.
In the opposite direction, the RNA sequence with 705 RT partners occurs across \num{25748} genomes
but only 17 species, and its partners range from 58 to \num{1254} amino acids within a single
family, indicating length and truncation variants of related proteins in a small number of heavily
sequenced taxa. Exact-sequence identity is a fine equivalence: a single residue difference, or a
partial gene call, produces a distinct sequence hash and therefore a distinct node.

The recurrence of a pair was classified by the kind of repetition it represents, so that database
redundancy is not mistaken for independent observation. Of the \num{30924} pairs, \num{12005} occur
once, \num{9362} recur only as multiple database copies of a single physical locus, \num{4452} recur
across multiple genomes of one species and \num{5056} across multiple species. Only the last two
classes support an evolutionary reading, and approximately one third of the apparent recurrence in
this resource is therefore database redundancy.

A one-to-one association as observed is useful for selecting a downstream dataset, but it does not
demonstrate orthogonality or exclusive biological pairing.

\subsection{Taxonomic representation and the effect of redundancy correction}
\label{sec:dbchar:taxonomy}

Genome identifiers resolved into their source catalogues essentially completely
(\SI{99.9999}{\percent} for \texttt{ncbi\_bacteria}, \SI{100}{\percent} for the remainder). A
resolved join is not, however, a resolved value: the NCBI assembly summaries carry no
genome-completeness column, so a completeness-based quality filter can speak for approximately
three-quarters of the corpus and is silent on the remainder.

Taxonomic rank coverage is a property of the schema rather than of data quality. Phylum is present
for \SI{100}{\percent} of records assigned under GTDB and for \SI{0}{\percent} of those assigned
under the NCBI scheme, by construction. Taxonomy is therefore reported per system and never pooled,
which matters in practice because the two schemes use different names for the same phyla, so that
pooling would silently divide single taxa across multiple rows. The corpus is overwhelmingly
bacterial, and archaeal records constitute a small minority.

Correcting for redundancy moves the distribution substantially (Table~\ref{tab:redundancy}).
Changing the unit from records to distinct RT proteins reduces the share of the most frequent NCBI
species from \SI{19.97}{\percent} to \SI{6.09}{\percent}, and the share of the ten most frequent
species from \SI{66.21}{\percent} to \SI{18.95}{\percent}. The former figure describes deposition
practice; the latter describes diversity.

Record counts measure sequencing effort. Any claim of over- or under-representation requires a
denominator of genomes surveyed, which a corpus assembled by the presence of an RT does not contain;
representation is reported here, and enrichment is not.

\subsection{Dependence on the annotation route}
\label{sec:dbchar:tools}

Across \num{663308} distinct Retron records, the three detection routes recovered markedly different
subsets: myRT \SI{92.77}{\percent}, PADLOC \SI{68.37}{\percent} and DefenseFinder
\SI{67.71}{\percent}, with all three agreeing on \SI{53.17}{\percent} (Figure~\ref{fig:venn}). Where
both subtype-writing tools produced a label (\num{365708} records), the labels agree after
normalisation on \num{160381} records (\SI{43.86}{\percent}).

Most consequentially for downstream use, RNA carriage depends strongly on which tools called the
locus: \SI{89.23}{\percent} for myRT and PADLOC jointly, \SI{70.39}{\percent} for all three, and
\SI{13.50}{\percent} for myRT alone. A tool-defined subset of this corpus is therefore not a random
subset of it, and any rate computed on one is conditioned on the detector combination that produced
it.

Agreement between these tools does not constitute independent corroboration, since they share model
lineage. Further, the absence of a tool from a record cannot be distinguished from that tool never
having been run on the corresponding genome, so these percentages bound tool coverage from below.

\subsection{Construction of the RT--ncRNA association dataset}
\label{sec:dbchar:dataset}

The characterisation delivers a set of named datasets, each with an explicit inclusion rule and an
explicit statement of the claims it does not support (Table~\ref{tab:inventory}). The separation
between sequence-level and context-level populations is material: \num{16008} distinct RT proteins
are usable as sequences but lack defensible genomic context, and are retained in the former and
excluded from the latter.

For the RT--RNA association resource specifically, the architecture described in
Section~\ref{sec:dbchar:geometry} suggests a selection rule based on family, strand, direction,
adjacency and separation. Applied cumulatively to canonical placements, restriction to the Retron
family, to same-strand placements, to upstream placements, to placements with no intervening coding
sequence, and to a separation of at most 200~bp yields \num{190028} placements, corresponding to
\num{22526} distinct sequence pairs, \num{21610} distinct RT proteins, \num{11514} distinct RNA
sequences and \num{156741} physical loci, retaining \SI{74}{\percent} of canonical pairs
(Table~\ref{tab:cascade}). The strand and adjacency criteria are nearly without cost because the
population already satisfies them; the separation threshold and the treatment of truncated windows
are the only consequential choices. Additionally excluding placements whose window is truncated at
the contig start reduces the resource to \num{14905} pairs, a further reduction of
\SI{34}{\percent}.

Within the \num{22526}-pair rule, \num{8660} pairs occur once, \num{6686} recur only as database
copies of a single physical locus, and \num{7180} recur across genomes or species.

The resource is therefore released in three tiers---permissive, strict, and
independence-filtered---rather than as a single thresholded set, so that a downstream application
may select its own operating point. Two limitations apply to all tiers. First, covariance-model hit
quality is not used: the score and expectation value accompanying every call enter no criterion, and
no tier should be described as high-confidence until they do. Second, exact-sequence identity
over-counts distinctness on both sides of the association, so sequences should be clustered at a
stated identity threshold before multiplicity is interpreted or any train/test partition is
constructed. Partitioning on unclustered sequence hashes would place boundary variants of a single
RNA, and truncation variants of a single protein, on both sides of the split.

Throughout, a canonical placement denotes one that is technically eligible and de-duplicated under
the stated geometric rules. It does not denote a biologically validated pair.

% =====================================================================
%  TABLES
% =====================================================================

\begin{table}[htbp]
\centering
\caption[Anchor populations of the mining corpus]{Anchor populations of the mining corpus. The
ncRNA-anchored records are excluded from every RT denominator reported in this chapter.}
\label{tab:populations}
\begin{tabular}{llrr}
\toprule
Population & Definition & Records & Files \\
\midrule
RT, single family & RT-anchored, one family label   & \num{3050688} & 41 \\
RT, multi-label   & RT-anchored, several labels     & \num{9012}    & 1  \\
ncRNA-anchored    & anchored on a non-coding RNA    & \num{298482}  & 1  \\
\midrule
Total             &                                 & \num{3358182} & 43 \\
\bottomrule
\end{tabular}
\end{table}

\begin{table}[htbp]
\centering
\caption[Analytical units of the corpus]{Analytical units of the corpus. The units are not
interchangeable, and the conversion between them varies by source collection.}
\label{tab:units}
\begin{tabular}{llr}
\toprule
Unit & Definition & Count \\
\midrule
Raw record              & one line of an RT-anchored file                      & \num{3059700} \\
Distinct record         & after byte-identical duplicates collapse             & \num{3051238} \\
Locus                   & distinct contig accession, start, end, strand        & \num{2847312} \\
Physical locus          & locus after contig-prefix normalisation              & \num{2475684} \\
Distinct RT protein     & distinct RT amino-acid sequence                      & \num{501561}  \\
Genome                  & distinct normalised genome identifier                & \num{1542433} \\
RT taxonomic occurrence & distinct protein--genome combination                 & \num{2737189} \\
\bottomrule
\end{tabular}
\end{table}

\begin{table}[htbp]
\centering
\caption[Classification of records lacking an annotated RT coding sequence]{Classification of the
\num{31504} RT-anchored records lacking an annotated RT coding sequence. The sequence-only and
ill-posed categories derive entirely from metagenome-assembled genomes.}
\label{tab:recovery}
\begin{tabular}{llr}
\toprule
State & Interpretation & Records \\
\midrule
Recovered     & coding sequence reconstructed; context retained & \num{14188} \\
Sequence only & protein usable; no defensible genomic context   & \num{16688} \\
Ill-posed     & neither coordinates nor context recoverable     & 628 \\
\midrule
\multicolumn{2}{l}{\emph{Sequence-only, by representation class}} & \\
\quad RT beyond the retrieved contig      & & \num{9128} \\
\quad RT crossing a contig boundary       & & \num{3636} \\
\quad Inverted window, no sequence        & & \num{3924} \\
\bottomrule
\end{tabular}
\end{table}

\begin{table}[htbp]
\centering
\caption[Geometry of canonical RT--ncRNA placements]{Geometry of the \num{344154} canonical
RT--ncRNA placements. All quantities are derived from coordinates and strand.}
\label{tab:geometry}
\begin{tabular}{lrr}
\toprule
Property & Placements & Percentage \\
\midrule
RNA upstream of the RT    & \num{325269} & 94.51 \\
Overlapping               & \num{12124}  & 3.52 \\
RNA downstream of the RT  & \num{6761}   & 1.96 \\
\midrule
Same strand as the RT     & \num{341113} & 99.12 \\
No intervening CDS        & \num{324904} & 94.41 \\
One intervening CDS       & \num{12427}  & 3.61 \\
More than three           & \num{5340}   & 1.55 \\
\midrule
\multicolumn{3}{l}{Median upstream separation: $-55$~bp} \\
\bottomrule
\end{tabular}
\end{table}

\begin{table}[htbp]
\centering
\caption[Atypical placements retained]{Atypical placements retained rather than filtered, on the
eligible placement population ($n=\num{345313}$).}
\label{tab:atypical}
\begin{tabular}{lr}
\toprule
Class & Placements \\
\midrule
Overlapping the RT                      & \num{12259} \\
More than two intervening CDS           & \num{5903} \\
Opposite strand                         & \num{3095} \\
Separated by more than 5~kb             & 557 \\
Locus carrying several distinct RNAs    & 398 \\
\bottomrule
\end{tabular}
\end{table}

\begin{table}[htbp]
\centering
\caption[Effect of redundancy correction on taxonomic composition]{Effect of redundancy correction
on taxonomic composition, reported separately for each taxonomy system.}
\label{tab:redundancy}
\begin{tabular}{llrr}
\toprule
System & Measure & On records (\%) & On distinct proteins (\%) \\
\midrule
NCBI & most frequent species    & 19.97 & 6.09 \\
NCBI & ten most frequent        & 66.21 & 18.95 \\
GTDB & most frequent species    & 8.57  & 2.21 \\
GTDB & ten most frequent        & 26.86 & 9.85 \\
\bottomrule
\end{tabular}
\end{table}

\begin{table}[htbp]
\centering
\caption[Cumulative construction of the RT--ncRNA association resource]{Cumulative construction of
the RT--ncRNA association resource from canonical placements.}
\label{tab:cascade}
\begin{tabular}{lrrrr}
\toprule
Criterion applied & Placements & Pairs & RT proteins & RNA sequences \\
\midrule
Canonical placements        & \num{344154} & \num{30427} & \num{28976} & \num{16011} \\
\quad + Retron family only  & \num{343892} & \num{30287} & \num{28838} & \num{15906} \\
\quad + same strand         & \num{340961} & \num{29837} & \num{28465} & \num{15666} \\
\quad + upstream            & \num{322223} & \num{26293} & \num{25043} & \num{13793} \\
\quad + no intervening CDS  & \num{311279} & \num{25295} & \num{24134} & \num{13252} \\
\quad + within 200~bp       & \num{190028} & \num{22526} & \num{21610} & \num{11514} \\
\quad + window not truncated& \num{154066} & \num{14905} & \num{14268} & \num{8382}  \\
\bottomrule
\end{tabular}
\end{table}

% Table:inventory is large; see the figure/table manifest for its column
% specification. Generated from tables/K1_dataset_inventory.tsv.
\begin{table}[htbp]
\centering
\caption[Dataset inventory]{Inventory of the datasets delivered by the characterisation, with
inclusion rule, size and unsupported claims. \emph{Placeholder: generate from
\texttt{K1\_dataset\_inventory.tsv}.}}
\label{tab:inventory}
\begin{tabular}{p{0.28\textwidth}p{0.30\textwidth}rp{0.24\textwidth}}
\toprule
Dataset & Inclusion rule & Size & Not usable for \\
\midrule
\multicolumn{4}{c}{\emph{[generate rows]}} \\
\bottomrule
\end{tabular}
\end{table}

% =====================================================================
%  FIGURES
% =====================================================================

\begin{figure}[htbp]
\centering
\includegraphics[width=\textwidth]{A1_unit_ladder}
\caption[Analytical unit ladder]{Reduction of the corpus across analytical units, and the
dependence of the locus-to-protein ratio on the source collection.}
\label{fig:units}
\end{figure}

\begin{figure}[htbp]
\centering
\includegraphics[width=\textwidth]{A4_database_attribution}
\caption[Cross-collection overlap]{Overlap between source collections. Per-collection counts are
coverage, not shares; distinct RT proteins cannot be attributed to a single collection.}
\label{fig:dboverlap}
\end{figure}

\begin{figure}[htbp]
\centering
\includegraphics[width=\textwidth]{B1_family_composition}
\caption[RT family composition]{Family composition of distinct RT proteins, compared with the
record-level composition.}
\label{fig:families}
\end{figure}

\begin{figure}[htbp]
\centering
\includegraphics[width=0.85\textwidth]{B2_rt_length_by_family}
\caption[RT protein length by family]{Distribution of RT protein length by family, on distinct
single-family proteins.}
\label{fig:lengths}
\end{figure}

\begin{figure}[htbp]
\centering
\includegraphics[width=0.85\textwidth]{B3_completeness_by_family}
\caption[Completeness composition by family]{Completeness composition by family. A missing Prodigal
attribute is classified as no evidence rather than as complete.}
\label{fig:completeness}
\end{figure}

\begin{figure}[htbp]
\centering
\includegraphics[width=0.8\textwidth]{FIG_multi_margin}
\caption[Profile score margins for multi-labelled proteins]{Margin between the best- and
second-best-scoring family profile for multi-labelled proteins, against a seeded single-family
control. \emph{Placeholder: figure not yet generated.}}
\label{fig:multi}
\end{figure}

\begin{figure}[htbp]
\centering
\includegraphics[width=0.85\textwidth]{C2_ncrna_length_by_model}
\caption[ncRNA length by covariance model]{Length distribution of detected RNAs by covariance
model. These are model-defined hit boundaries, not transcript boundaries.}
\label{fig:ncrnalen}
\end{figure}

\begin{figure}[htbp]
\centering
\includegraphics[width=0.85\textwidth]{C3_model_composition}
\caption[Covariance-model composition]{Model composition at the level of calls and of distinct RNA
sequences; the difference measures the redundancy of each model's matches.}
\label{fig:cmcomp}
\end{figure}

\begin{figure}[htbp]
\centering
\includegraphics[width=\textwidth]{D2_distance_distribution}
\caption[RT--ncRNA separation]{Signed separation between RT and RNA, stratified by truncation of the
window at the contig start, which isolates the technical downstream mode.}
\label{fig:geometry}
\end{figure}

\begin{figure}[htbp]
\centering
\includegraphics[width=\textwidth]{D6_joint_geometry}
\caption[Joint geometry and weighting]{Joint distribution of direction, strand, adjacency and
separation, and the effect of re-weighting from placements to physical loci and to distinct
sequence pairs.}
\label{fig:joint}
\end{figure}

\begin{figure}[htbp]
\centering
\includegraphics[width=\textwidth]{D5_strand_agreement}
\caption[Strand agreement and the opposite-strand population]{Same-strand rate across populations,
and the structure of the opposite-strand placements.}
\label{fig:strand}
\end{figure}

\begin{figure}[htbp]
\centering
\includegraphics[width=\textwidth]{F_pairing_topology}
\caption[Pairing topology and recurrence]{Partner-count distributions on both sides of the
association, component shapes, and the classification of pair recurrence.}
\label{fig:topology}
\end{figure}

\begin{figure}[htbp]
\centering
\includegraphics[width=0.8\textwidth]{FIG_tool_venn}
\caption[Agreement between detection routes]{Agreement between myRT, PADLOC and DefenseFinder on
distinct Retron RT proteins, annotated with RNA carriage per intersection. \emph{Placeholder: figure
not yet generated.}}
\label{fig:venn}
\end{figure}
